%% ---------------------------------------------------------------------
\section{Preliminaries}\label{sec:prelim}

\subsection{The problem}

Let $S$ be a finite set of \emph{states} containing three distinguished
elements: the \emph{quiescent} state $\mathsf L$, the \emph{general}
$\mathsf G$ and the \emph{firing} state $\mathsf F$; these are pairwise
distinct. We write $W:=S\setminus\{\mathsf F\}$ for the set of
\emph{working} states and use a symbol $\ast\notin S$ for the border.
A \emph{rule} is a map
\[
\delta:\;(W\cup\{\ast\})\times W\times(W\cup\{\ast\})\;\longrightarrow\;S
\]
satisfying the quiescence conditions
\begin{equation}\label{eq:quiescence}
\delta(\mathsf L,\mathsf L,\mathsf L)=\delta(\mathsf L,\mathsf L,\ast)=\mathsf L .
\end{equation}
(Transitions out of $\mathsf F$ and transitions whose neighbourhood contains
$\mathsf F$ play no role below and are left unspecified; we do \emph{not}
require $\delta(\ast,\mathsf L,\mathsf L)=\mathsf L$, which only strengthens
the non-existence statements.)

For $n\ge 1$ the \emph{line of length $n$} is the space-time diagram
$C_n:\mathbb N\times\{1,\dots,n\}\to S$ defined by
$C_n(0,1)=\mathsf G$, $C_n(0,i)=\mathsf L$ for $2\le i\le n$, and
\[
C_n(t+1,i)=\delta\bigl(C_n(t,i-1),\,C_n(t,i),\,C_n(t,i+1)\bigr),
\]
with the convention $C_n(t,0)=C_n(t,n+1)=\ast$; the diagram is only
considered up to the first time at which some cell is in state $\mathsf F$.
The rule \emph{synchronizes the line of length $n$ at time $T$} if
$C_n(T,i)=\mathsf F$ for all $1\le i\le n$ and $C_n(t,i)\neq\mathsf F$ for all
$t<T$ and all~$i$.

It is classical \cite{Moore64,Waksman66,Balzer67} that no rule synchronizes a
line of length $n\ge2$ before time $2n-2$ (see also \cref{lem:return}
below). A \emph{minimal-time solution} is a rule that synchronizes the line
of length $n$ at time exactly $2n-2$ for every $n\ge2$. (For $n=1$ one
usually asks $\delta(\ast,\mathsf G,\ast)=\mathsf F$; this transition is
used by no other line and is irrelevant here.) A \emph{$k$-state solution} is
one with $|S|=k$.

More generally, for a set $\mathcal N\subseteq\{2,3,\dots\}$ of lengths, a
rule is \emph{$\mathcal N$-minimal} if it synchronizes the line of length $n$
at time $2n-2$ for every $n\in\mathcal N$. Clearly a minimal-time solution
is $\mathcal N$-minimal for every $\mathcal N$; all our non-existence
arguments exhibit a finite family of necessary conditions that no $5$-state
rule satisfies.

\subsection{The half-line and the agreement lemma}

Let $C_\infty:\mathbb N\times\{1,2,\dots\}\to S$ be the diagram of the
half-line: $C_\infty(0,1)=\mathsf G$, $C_\infty(0,i)=\mathsf L$ for $i\ge2$,
and the same recurrence with a border on the left only.

\begin{lemma}[light cone]\label{lem:lightcone}
For every rule, every $n$ and every $(t,i)$ with $t<i-1$ we have
$C_n(t,i)=C_\infty(t,i)=\mathsf L$.
\end{lemma}

\begin{proof}
By induction on $t$. For $t=0$ this is the initial condition. Let $t\ge1$
and $t<i-1$. The three cells $(t-1,i-1),(t-1,i),(t-1,i+1)$ satisfy
$t-1<(i-1)-1$, hence are in state $\mathsf L$ by induction, except that the
right neighbour of cell $n$ of $C_n$ is the border. By \eqref{eq:quiescence}
the new state is $\mathsf L$ in both cases.
\end{proof}

\begin{lemma}[agreement]\label{lem:agree}
For every rule and every $n\ge 2$: $C_n(t,i)=C_\infty(t,i)$ whenever
$1\le i\le n$ and $t+i\le 2n-2$.
\end{lemma}

\begin{proof}
Induction on $t$; the case $t=0$ is clear. Let $t\ge1$ and $t+i\le 2n-2$.
If $i=n$ then $t\le n-2<i-1$ and both sides equal $\mathsf L$ by
\cref{lem:lightcone}. If $i<n$, the three neighbours $(t-1,j)$,
$j\in\{i-1,i,i+1\}$, satisfy $(t-1)+j\le t+i\le 2n-2$ and $j\le n$, so they
agree by induction (for $j=0$ both are the border). Hence the new states
agree.
\end{proof}

Thus the right border of the line of length $n$ is first ``felt'' at the
cell $(n-1,n)$, and its influence is confined to the \emph{reflected
triangle}
\[
R_n:=\{(t,i):\ 1\le i\le n,\ t+i\ge 2n-1,\ t\le 2n-2\}.
\]

\begin{lemma}[return chain]\label{lem:return}
Let $\delta$ synchronize the lines of lengths $n$ and $n+1$ at times $2n-2$
and $2n$ respectively. Then
\[
C_n(2n-1-i,\,i)\;\neq\;C_\infty(2n-1-i,\,i)\qquad(1\le i\le n).
\]
In particular the \emph{front} $s_i:=C_\infty(i-1,i)$ satisfies
$s_{n-1}\neq\mathsf L$, and $\delta(s_{n-1},\mathsf L,\ast)\neq
\delta(s_{n-1},\mathsf L,\mathsf L)$.
\end{lemma}

\begin{proof}
For $i=1$: $C_n(2n-2,1)=\mathsf F$ while, by \cref{lem:agree} applied to
the line of length $n+1$ (note $(2n-2)+1\le 2(n+1)-2$),
$C_\infty(2n-2,1)=C_{n+1}(2n-2,1)\neq\mathsf F$ because $2n-2<2n$.
Suppose the claim holds for some $i<n$. The cell $(2n-1-i,i)$ is computed
from the cells $(2n-2-i,j)$, $j=i-1,i,i+1$, whose coordinate sums are
$2n-3$, $2n-2$ and $2n-1$. By \cref{lem:agree} the first two are the same
in $C_n$ and $C_\infty$ (for $i=1$ the left neighbour is the border in both
diagrams). Since the results differ, the right neighbours differ:
$C_n(2n-2-i,i+1)\ne C_\infty(2n-2-i,i+1)$, which is the claim for $i+1$.
For $i=n$ we get $C_n(n-1,n)=\delta(s_{n-1},\mathsf L,\ast)\neq
C_\infty(n-1,n)=\delta(s_{n-1},\mathsf L,\mathsf L)=s_n$, which is impossible
if $s_{n-1}=\mathsf L$ by \eqref{eq:quiescence}.
\end{proof}

Applying \cref{lem:return} to all $n$ shows that in a minimal-time
solution every cell $i$ leaves $\mathsf L$ exactly at time $i-1$ (the front
moves at speed one) and, since $C_\infty(t,i)=C_{n}(t,i)$ for $n$ large,
\begin{equation}\label{eq:neverfire}
C_\infty(t,i)\neq\mathsf F\qquad\text{for all }t\ge0,\ i\ge1 .
\end{equation}
